\section*{Chapter \ref{ch:pl:a-tick-store-on-open-formats} --- A Tick Store on Open Formats}

\begin{solution}{exo:pl:a-tick-store-on-open-formats:1}
$14\,478\,152 - 258\,535 \times 56 = 192$ bytes of header. The version-2 header of the week is 256 bytes (its layout, as JSON,
is longer), so the file would be $256 + 258\,535 \times 64 = 16\,546\,496$ bytes.
\end{solution}

\begin{solution}{exo:pl:a-tick-store-on-open-formats:2}
Bytes already written never change, and records have a fixed width, so the first
$\lfloor (\text{size} - \text{header}) / \text{width} \rfloor$ records are complete and final. The reader ignores the partial
record at the tail (the writer is in the middle of it) and reads it on its next pass.
\end{solution}

\begin{solution}{exo:pl:a-tick-store-on-open-formats:3}
$17.5 / 6.15 = 2.85$. The partitions are sorted by symbol and time and stored as columns, so each column is encoded (dictionaries,
deltas) and compressed; the flat records keep every field at full width, and the IPC streams are uncompressed.
\end{solution}

\begin{solution}{exo:pl:a-tick-store-on-open-formats:4}
223 of the 17\,281 trades, 1.3\%: for each, the join returns the quote the execution produced, with the executed level removed.
The mid has already moved in the trade's direction, so the trade looks closer to the mid than it was: the measured effective
spread is biased toward zero, and a trade sign inferred from that quote can flip and make the spread negative.
\end{solution}

\begin{solution}{exo:pl:a-tick-store-on-open-formats:5}
About 2.2 million events a second for 35 minutes, 2\,100 seconds: about 4.7 billion events on one core. The options feed's
planned 311 billion messages a day is over sixty times more: the compaction must run in parallel by partition on
over sixty cores, or incrementally during the day (hour by hour), or both.
\end{solution}

\begin{solution}{exo:pl:a-tick-store-on-open-formats:6}
If the version number is not changed, an old reader reads the new field's bytes as the fields that follow it: silently wrong
values, the worst failure. If it is changed and the width grows, an old reader refuses the file. \texttt{check\_evolution} reports
the new field as inserted before existing fields.
\end{solution}

\begin{solution}{exo:pl:a-tick-store-on-open-formats:7}
Instrument~1: 27\,157 records; instrument~2: 102\,144 records, 367\,600 shares executed; instrument~0 (system messages): 4
records. The C++ and Rust readers give the same numbers on the fixtures, which are written by the same Python code.
\end{solution}

\begin{solution}{exo:pl:a-tick-store-on-open-formats:8}
A crash during the compaction leaves the day neither in its intraday form nor complete in its historical form, and a query during
it reads a mixture. Compact to new files, verify them (row counts, sequence ranges), then switch readers to them atomically (a
rename of the partition directory or a manifest), and delete the intraday files only after that.
\end{solution}

\begin{solution}{pb:pl:a-tick-store-on-open-formats:1}
\begin{enumerate}
\item 564\,162 events, 96\,206 quotes and 17\,281 trades.
\item The magic number, the version, the record width and the layout as JSON; padded to 64 bytes so that every record starts at
the same alignment in every file.
\item The flat events are the layout of the trading path, for replay and C++ or Rust tools; the IPC tables carry a schema and
strings, for dataframes and SQL.
\item 56 and 64 bytes.
\item The session must write at the rate data arrives; sorting needs the whole day, and compression costs time the capture does not
have.
\item Partitioned by date and symbol, sorted by time inside: queries name a date range and a symbol and filter a time window.
\item 17.5~MB intraday, 6.15~MB compacted.
\item 0.12~s for 258\,535 events, about 2.2 million events a second on one core.
\item About 4.7 billion.
\item Compact partitions in parallel, compact incrementally during the session, and keep the intraday form cheap to convert.
\item The query layer defines one view per table as the union, by column name, of the historical partitions and the current day's
table.
\item 12\,291, 996, 1\,388, 1\,197 and 4\,473.
\item An execution and the quote change it causes share a timestamp; the sequence number orders them.
\item 223 trades; DuckDB, Polars and the reference agree on every row with the sequence number.
\item 16~ms, 45~ms and 0.17~s.
\item \textbf{Named result.} Monday's 258\,535 events compact in 0.12~s (about 2.2 million events a second on one core) from
17.5~MB to 6.15~MB; a five-day window query takes 18~ms and the week's as-of join 16~ms in DuckDB; at that rate the 35-minute
window absorbs about 4.7 billion events, about one sixty-sixth of the options feed's planned day.
\item The old days have no venue column and read it as null; the flat files carry their version and are upgraded on read with the
field at its default.
\item A scan four times faster (1.2 against 4.8~ms) with no decoding, at 2.8 times the size.
\item The intraday/historical split and the end-of-day write to history come from the tick databases; the formats, the query engine
and the dataframes are open.
\item One query layer that matches columns by name over the current day's tables and the historical partitions.
\end{enumerate}
\end{solution}

\begin{solution}{iq:pl:a-tick-store-on-open-formats:1}
They want opposite layouts: the current day must be written as it arrives (appended, in time order, unsorted), history must be
read in ranges across days and symbols (partitioned, sorted, compressed). Compaction converts one into the other once, at the
close, and a query layer hides the split.
\iqlookfor{Write-optimised against read-optimised, and the compaction between them.}
\end{solution}

\begin{solution}{iq:pl:a-tick-store-on-open-formats:2}
The as-of join on the timestamp can return the quote the trade itself produced, since both carry the same exchange time: the mid
has already moved and the spread looks smaller or negative. Join strictly on the feed's sequence number (the last quote of an
earlier sequence number), or on a receive-order key.
\iqlookfor{Ties at equal timestamps; strict join on a total order.}
\end{solution}

\begin{solution}{iq:pl:a-tick-store-on-open-formats:3}
Reads without system calls or copies: pages are loaded on first touch and shared between processes, and fixed-width records are
addressed directly. It is a bad idea for files larger than memory with random access (page faults dominate), for files that can be
truncated under the reader, and where latency must be bounded (a page fault on the hot path).
\iqlookfor{Page faults, sharing, and the failure cases.}
\end{solution}

\begin{solution}{iq:pl:a-tick-store-on-open-formats:4}
Add the field at the end with a default, raise the version in the headers and the registry, make every reader accept every
version, and leave old files as they are: old readers ignore the field, new readers see the default for old days.
\iqlookfor{Add-only evolution with versions; no rewrite of history.}
\end{solution}

\begin{solution}{iq:pl:a-tick-store-on-open-formats:5}
Partition by date (every query has a date range) and, if files stay large enough, by symbol or a hash bucket of symbols; sort by
symbol then time inside; \omterm{def:pl:columnar-formats:rowgroup}{row groups} sized by measurement. 8\,000 symbols a day would make tiny files per symbol for the illiquid
ones, so bucket them.
\iqlookfor{Date partitions, symbol order, and the small-files problem.}
\end{solution}

\begin{solution}{iq:pl:a-tick-store-on-open-formats:6}
Capture to append-only flat files and Arrow IPC during the session; compaction at the close to Parquet partitioned by date and
symbol, sorted by time, written to new files and switched atomically; DuckDB or Polars over both through views matched by name;
versioned schemas. Measure first the bytes per event in each form, the compaction rate against the close-to-research window, and
the latency of the main queries (as-of joins, windows).
\iqlookfor{The whole cycle with an atomic switch, and the measurements that size it.}
\end{solution}
